\documentclass{article}

\usepackage{fancyhdr}
\usepackage{extramarks}
\usepackage{amsmath}
\usepackage{amsthm}
\usepackage{amsfonts}
\usepackage{tikz}
\usepackage[plain]{algorithm}
\usepackage{algpseudocode}

\usetikzlibrary{automata,positioning}

%
% Basic Document Settings
%

\topmargin=-0.45in
\evensidemargin=0in
\oddsidemargin=0in
\textwidth=6.5in
\textheight=9.0in
\headsep=0.25in

\linespread{1.1}

\pagestyle{fancy}
\lhead{\hmwkAuthorName}
\chead{\hmwkClass , \hmwkClassInstructor: \hmwkTitle}
\rhead{\firstxmark}
\lfoot{\lastxmark}
\cfoot{\thepage}

\renewcommand\headrulewidth{0.4pt}
\renewcommand\footrulewidth{0.4pt}

\setlength\parindent{0pt}

%
% Create Problem Sections
%

\newcommand{\enterProblemHeader}[1]{
    \nobreak\extramarks{}{Problem \arabic{#1} continued on next page\ldots}\nobreak{}
    \nobreak\extramarks{Problem \arabic{#1} (continued)}{Problem \arabic{#1} continued on next page\ldots}\nobreak{}
}

\newcommand{\exitProblemHeader}[1]{
    \nobreak\extramarks{Problem \arabic{#1} (continued)}{Problem \arabic{#1} continued on next page\ldots}\nobreak{}
    \stepcounter{#1}
    \nobreak\extramarks{Problem \arabic{#1}}{}\nobreak{}
}

\setcounter{secnumdepth}{0}
\newcounter{partCounter}
\newcounter{homeworkProblemCounter}
\setcounter{homeworkProblemCounter}{1}
\nobreak\extramarks{Problem \arabic{homeworkProblemCounter}}{}\nobreak{}

%
% Homework Problem Environment
%
% This environment takes an optional argument. When given, it will adjust the
% problem counter. This is useful for when the problems given for your
% assignment aren't sequential. See the last 3 problems of this template for an
% example.
%
\newenvironment{homeworkProblem}[1][-1]{
    \ifnum#1>0
        \setcounter{homeworkProblemCounter}{#1}
    \fi
    \section{Problem \arabic{homeworkProblemCounter}}
    \setcounter{partCounter}{1}
    \enterProblemHeader{homeworkProblemCounter}
}{
    \exitProblemHeader{homeworkProblemCounter}
}

%
% Homework Details
%   - Title
%   - Due date
%   - Class
%   - Section/Time
%   - Instructor
%   - Author
%

\newcommand{\hmwkTitle}{Problem Set 1}
\newcommand{\hmwkDueDate}{DUE DATE}
\newcommand{\hmwkClass}{CS229}
\newcommand{\hmwkClassTime}{}
\newcommand{\hmwkClassInstructor}{Prof. Andrew Ng}
\newcommand{\hmwkAuthorName}{\textbf{Connor Yan}}

%
% Title Page
%

\title{
    \vspace{2in}
    \textmd{\textbf{\hmwkClass}}\\
    \textmd{\textbf{\hmwkTitle}}\\
    \vspace{3in}
}

\author{\hmwkAuthorName}
\date{}

\renewcommand{\part}[1]{\textbf{\large Part \Alph{partCounter}}\stepcounter{partCounter}\\}

%
% Various Helper Commands
%

% Useful for algorithms
\newcommand{\alg}[1]{\textsc{\bfseries \footnotesize #1}}

% For derivatives
\newcommand{\deriv}[1]{\frac{\mathrm{d}}{\mathrm{d}x} (#1)}

% For partial derivatives
\newcommand{\pderiv}[2]{\frac{\partial}{\partial #1} (#2)}

% Integral dx
\newcommand{\dx}{\mathrm{d}x}

% Alias for the Solution section header
\newcommand{\solution}{\textbf{\large Solution}}

% Probability commands: Expectation, Variance, Covariance, Bias
\newcommand{\E}{\mathrm{E}}
\newcommand{\Var}{\mathrm{Var}}
\newcommand{\Cov}{\mathrm{Cov}}
\newcommand{\Bias}{\mathrm{Bias}}
\newcommand{\R}{\mathbb{R}}
\newcommand{\mb}{\mathrm}

\begin{document}

\maketitle

\pagebreak

\begin{homeworkProblem}

    \textbf{(a)}
    The average empirical loss for logistic regression:
    \[
        J(\theta)=-\frac{1}{m}\sum^m_{i=1}
        {
            y^{(i)}\log{\left(h_\theta (x^{(i)})\right)+
            (1-y^{(i)})\log{\left(1-h_\theta(x^{(i)})\right)}}
        }
    \]
    where \$y^{(i)}\in\{0,1\}$, $h_\theta(x)=g(\theta^Tx)$ and $g(z)=1/(1+e^{-z})$. As $g'(z)=g(z)(1-g(z))$, therefore the gradient is
    \[
    \begin{aligned}
        \nabla_\theta J(\theta)
        &=-\frac{1}{m}\sum^m_{i=1}
        {
            y^{(i)}\frac{h_\theta (x^{(i)})(1-h_\theta (x^{(i)}))x^{(i)}}{h_\theta (x^{(i)})}
            +(1-y^{(i)})\frac{-h_\theta (x^{(i)})(1-h_\theta (x^{(i)}))x^{(i)}}{\left(1-h_\theta(x^{(i)})\right)}
        }\\
        &=-\frac{1}{m}\sum^m_{i=1}
        {
            y^{(i)}(1-h_\theta (x^{(i)}))x^{(i)}
            -(1-y^{(i)})h_\theta (x^{(i)})x^{(i)}
        }\\
        &=-\frac{1}{m}\sum^m_{i=1}
        {
            \left(y^{(i)}-h_\theta (x^{(i)})\right)x^{(i)}
        }
    \end{aligned}
    \]
     
     \[
     \begin{aligned}
        H
        &=-\frac{1}{m}\nabla_\theta\sum^m_{i=1}
        {
            \left(y^{(i)}-h_\theta (x^{(i)})\right)x^{(i)}
        }\\
        &=\frac{1}{m}\sum^m_{i=1}
        {
            \nabla_\theta\left(h_\theta (x^{(i)})x^{(i)}\right)
        }\\
        &=\frac{1}{m}\sum^m_{i=1}
        {
            (1-h_\theta (x^{(i)}))(h_\theta (x^{(i)})x^{(i)}(x^{(i)})^T
        }.
    \end{aligned}
    \]

    For any vector $z\in\R^n$, we want to prove
    \[
        z^T Hz\ge0.
    \]
    First take
    \[
    \begin{aligned}
        z^TH
        &=\begin{bmatrix}
            \sum^n_{i=1}z_iH_{i1} & \cdots & \sum^n_{i=1}z_nH_{in}
        \end{bmatrix}
    \end{aligned}
    \]
    Thus, 
    \[
    \begin{aligned}
        z^THz&=\sum^n_{j=1}{z_j\sum^n_{i=1}z_iH_{ij}}\\
        &=\sum^n_{i=1}{\sum^n_{j=1}{z_iz_jH_{ij}}}\\
        &=\frac{1}{m}\sum^m_{k=1}h_\theta(x^{(k)})(1-h_\theta({x^{(k)}}))\sum^n_{i=1}{\sum^n_{j=1}{z_iz_jx^{(k)}_ix^{(k)}_j}}
    \end{aligned}
    \]
    Note that
    \[
        \sum^n_{i=1}\sum^n_{j=1}{z_iz_jx^{(k)}_ix^{(k)}}
        =\left(\sum^n_{i=1}{z_ix^{(k)}_i}\right)\cdot\left(\sum^n_{j=1}{z_jx^{(k)}_j}\right)
        =(z^Tx^{(k)})^2\ge0
    \]
    And note that as $h_\theta\in[0,1]$, $h_\theta(x^{(k)})(1-h_\theta{x^{(k)}})$ is also greater or equal to zero. Hence, $z^THz\ge0$.
    
    \textbf{(b)} See code.

    \textbf{(c)}

    \[
    \begin{aligned} 
    p(y=1|x;\phi,\mu_0,\mu_1,\Sigma)
    &=\frac{p(x|y=1)p(y=1)}{p(x)} \\
    &=\frac{p(x|y=1)p(y=1)}{p(x|y=0) p(y=0)+p(x|y=1)p(y=1)} \\
    &= 1/\left(1+\frac{p(x|y=0)p(y=0)}{p(x|y=1)p(y=1)}\right) \\
    &=1/\left(1+\exp{\left(\frac{1-\phi}{\phi} \frac{1}{2}((x-\mu_1)^T\Sigma^{-1}(x-\mu_1)-(x-\mu_0)^T\Sigma^{-1}(x-\mu_0))\right)}\right) \\
    &=1/\left(1+ \exp{\left((\mu_0-\mu_1)^T\Sigma^{-1}x+\frac{1}{2}(\mu_1^T\Sigma^{-1}\mu_1-\mu_0^T\Sigma^{-1}\mu_0)+\log \frac{1-\phi}{\phi}\right)}\right) \\
    &=\frac{1}{1+\exp{(-(\theta^T x+\theta_0))}},
    \end{aligned}
    \]

    where $\theta^T=(\mu_0-\mu_1)^T\Sigma^{-1}$ and $\theta_0=-\frac{1}{2}(\mu_1^T\Sigma^{-1}\mu_1-\mu_0^T\Sigma^{-1}\mu_0)-\log \frac{1-\phi}{\phi}$.

    \textbf{(d)}
    \[
    \begin{aligned}
    \ell(\phi,\mu_0,\mu_1,\Sigma)
     &=\sum_{i=1}^m\log{\left(p(x^{(i)}|y^{(i)};\mu_0,\mu_1,\Sigma)\right)}+\log{\left(p(y^{(i)};\phi)\right)}\\
     &=\sum_{i=1}^m{
        -\frac{1}{2}(x^{(i)}-\mu_{y^{(i)}})^T\Sigma^{-1}(x^{(i)}-\mu_{y^{(i)}})
        -\log((2\pi)^{n/2}|\Sigma|^{1/2})
        +\log(\phi^{y^{(i)}}(1-\phi)^{1-y^{(i)}})
     }
    \end{aligned}
    \]
    w.r.t. $\phi$:

    \[
    \begin{aligned}
    \frac{\partial}{\partial\phi}\ell(\phi,\mu_0,\mu_1,\Sigma)
     &=\frac{\partial}{\partial\phi}\sum_{i=1}^m{
        \log{(\phi^{y^{(i)}})}+\log{((1-\phi)^{1-y^{(i)}})}
     } \\
     &=\sum_{i=1}^m{\frac{y^{(i)}}{\phi}}-\sum_{i=1}^m\frac{1-y^{(i)}}{1-\phi} =0
    \end{aligned}
    \]
    \[
    \begin{aligned}
    \frac{1}{\phi}\sum_{i=1}^m1\{y^{(i)}=1\}&=\frac{1}{1-\phi}\sum_{i=1}^m1\{y^{(i)}=0\}\\
    \left(\frac{1}{\phi}-1\right)\sum_{i=1}^m1\{y^{(i)}=1\}&=m-\sum_{i=1}^m1\{y^{(i)}=1\}\\
    \left(\frac{1}{\phi}-1\right)&=m/\sum_{i=1}^m1\{y^{(i)}=1\}-1\\
    \phi&=\frac{1}{m}\sum_{i=1}^m1\{y^{(i)}=1\}
    \end{aligned}
    \]
    w.r.t. $\mu_0$    \[
    \begin{aligned}
    \nabla_{\mu_0}\ell(\phi,\mu_0,\mu_1,\Sigma)
     &=\nabla_{\mu_0}\sum_{i=1}^m{
        -\frac{1\{y^{(i)}=0\}}{2}(x^{(i)}-\mu_0)^T\Sigma^{-1}(x^{(i)}-\mu_0)
     } \\
    &=\sum_{i=1}^m{
        1\{y^{(i)}=0\}(x^{(i)}-\mu_0)
     }=0
    \end{aligned}
    \]

    \[
    \begin{aligned}
    \sum_{i=1}^m{
        1\{y^{(i)}=0\}x^{(i)}
     }&=\mu_0\sum_{i=1}^m{
        1\{y^{(i)}=0\}
     } \\
     \mu_0&=\frac{\sum_{i=1}^m{1\{y^{(i)}=0\}x^{(i)}}}{\sum_{i=1}^m{1\{y^{(i)}=0\}}}
    \end{aligned}
    \]
    Similar for $\mu_1$, we have
    \[ \mu_1=\frac{\sum_{i=1}^m{1\{y^{(i)}=1\}x^{(i)}}}{\sum_{i=1}^m{1\{y^{(i)}=1\}}}\]
    w.r.t. $\Sigma$, consider $\Sigma$ as a scalar $\sigma^2$, hence:
    \[
    \begin{aligned}
        \frac{\partial}{\partial\Sigma}\ell(\phi,\mu_0,\mu_1,\Sigma)
        &= \frac{\partial}{\partial\Sigma}\sum_{i=1}^m{
        -\frac{1}{2}(x^{(i)}-\mu_{y^{(i)}})^T\Sigma^{-1}(x^{(i)}-    \mu_{y^{(i)}})
        -\frac{1}{2}\log|\Sigma|
        }\\
        &=-\sum_{i=1}^m{
        \frac{1}{2}(x^{(i)}-    \mu_{y^{(i)}})(x^{(i)}-\mu_{y^{(i)}})^T\Sigma^{-2}
        +\frac{1}{2\Sigma}
        }=0
    \end{aligned}
    \]
\[
\begin{aligned}
    -\sum_{i=1}^m{
        \frac{1}{2}(x^{(i)}-    \mu_{y^{(i)}})(x^{(i)}-\mu_{y^{(i)}})^T
        +\frac{m}{2}\Sigma
        }&=0\\
    \Sigma&=\frac{1}{m}\sum_{i=1}^m{
        (x^{(i)}-\mu_{y^{(i)}})
        (x^{(i)}-\mu_{y^{(i)}})^T
        }
\end{aligned}
\]
    

\end{homeworkProblem}

\newpage
\begin{homeworkProblem}
    
\end{homeworkProblem}

\end{document}
